% !TEX TS-program = XeLaTeX
% use the following command:
% all document files must be coded in UTF-8
\documentclass[portuguese]{textolivre}
% build HTML with: make4ht -e build.lua -c textolivre.cfg -x -u article "fn-in,svg,pic-align"

\journalname{Texto Livre}
\thevolume{19}
%\thenumber{1} % old template
\theyear{2026}
\receiveddate{\DTMdisplaydate{2025}{9}{22}{-1}} % YYYY MM DD
\accepteddate{\DTMdisplaydate{2026}{1}{6}{-1}}
\publisheddate{\DTMdisplaydate{2026}{6}{11}{-1}}
\corrauthor{João Rodrigo Maciel Portes}
\articledoi{10.1590/1983-3652.2026.61885}
%\articleid{NNNN} % if the article ID is not the last 5 numbers of its DOI, provide it using \articleid{} commmand 
% list of available sesscions in the journal: articles, dossier, reports, essays, reviews, interviews, editorial
\articlesessionname{articles}
\runningauthor{Portes et al.} 
%\editorname{Leonardo Araújo} % old template
\sectioneditorname{Daniervelin Pereira~\orcid{0000-0003-1861-3609}}
\layouteditorname{Leonardo Araujo~\orcid{0000-0003-3884-2177}}

\title{Avaliação da usabilidade de um aplicativo móvel para o treinamento parental no manejo da seletividade alimentar em crianças com transtorno do espectro autista}
\othertitle{Usability evaluation of a mobile application for parental training in the management of food selectivity in children with autism spectrum disorder}
% if there is a third language title, add here:
%\othertitle{Artikelvorlage zur Einreichung beim Texto Livre Journal}

\author[1]{João Rodrigo Maciel Portes~\orcid{0000-0003-2512-4491}\thanks{Email: \href{mailto:joaorodrigo@univali.br}{joaorodrigo@univali.br}}}
\author[2]{Weliton da Silva~\orcid{009-0008-3365-7865}\thanks{Email: \href{mailto:contatowelitonslv@gmail.com}{contatowelitonslv@gmail.com}}}
\author[2]{João Vitor Girardi Turqueti da Rocha~\orcid{0009-0006-0897-5613}\thanks{Email: \href{mailto:girardi.tr@gmail.com}{girardi.tr@gmail.com}}}
\author[1]{Letícia Corrêa Ferreira~\orcid{0000-0003-0739-0200}\thanks{Email: \href{mailto:ltccferreira@gmail.com}{ltccferreira@gmail.com}}}
\author[3]{Leila Bagaiolo~\orcid{0000-0002-5566-8230}\thanks{Email: \href{mailto:leila.bagaiolo@grupogradual.com.br}{leila.bagaiolo@grupogradual.com.br}}}
\author[3]{Claudia Pacífico~\orcid{0000-0002-9162-4614}\thanks{Email: \href{mailto:claudia.romano@grupogradual.com.br}{claudia.romano@grupogradual.com.br}}}
\author[4]{Mauro Luís Vieira~\orcid{0000-0003-0541-4133}\thanks{Email: \href{mailto:maurolvieira@gmail.com}{maurolvieira@gmail.com}}}
\affil[1]{Universidade do Vale do Itajaí, Mestrado Profissional de Psicologia, Itajaí, SC, Brasil.}
\affil[2]{Universidade do Vale do Itajaí, Curso de Psicologia, Itajaí, SC, Brasil.}
\affil[3]{Grupo de Intervenção Comportamental Gradual, São Paulo, SP, Brasil.}
\affil[4]{Universidade Federal de Santa Catarina, Programa de Pós-Graduação em Psicologia, Florianopólis, SC, Brasil.}

\addbibresource{article.bib}
% use biber instead of bibtex
% $ biber article

% used to create dummy text for the template file
\definecolor{dark-gray}{gray}{0.35} % color used to display dummy texts
\usepackage{lipsum}
\SetLipsumParListSurrounders{\colorlet{oldcolor}{.}\color{dark-gray}}{\color{oldcolor}}

% used here only to provide the XeLaTeX and BibTeX logos
\usepackage{hologo}

% if you use multirows in a table, include the multirow package
\usepackage{multirow}

% provides sidewaysfigure environment
\usepackage{rotating}

% CUSTOM EPIGRAPH - BEGIN 
%%% https://tex.stackexchange.com/questions/193178/specific-epigraph-style
\usepackage{epigraph}
\renewcommand\textflush{flushright}
\makeatletter
\newlength\epitextskip
\pretocmd{\@epitext}{\em}{}{}
\apptocmd{\@epitext}{\em}{}{}
\patchcmd{\epigraph}{\@epitext{#1}\\}{\@epitext{#1}\\[\epitextskip]}{}{}
\makeatother
\setlength\epigraphrule{0pt}
\setlength\epitextskip{0.5ex}
\setlength\epigraphwidth{.7\textwidth}
% CUSTOM EPIGRAPH - END

% to use IPA symbols in unicode add
%\usepackage{fontspec}
%\newfontfamily\ipafont{CMU Serif}
%\newcommand{\ipa}[1]{{\ipafont #1}}
% and in the text you may use the \ipa{...} command passing the symbols in unicode

% LANGUAGE - BEGIN
% ARABIC
% for languages that use special fonts, you must provide the typeface that will be used
% \setotherlanguage{arabic}
% \newfontfamily\arabicfont[Script=Arabic]{Amiri}
% \newfontfamily\arabicfontsf[Script=Arabic]{Amiri}
% \newfontfamily\arabicfonttt[Script=Arabic]{Amiri}
%
% in the article, to add arabic text use: \textlang{arabic}{ ... }
%
% RUSSIAN
% for russian text we also need to define fonts with support for Cyrillic script
% \usepackage{fontspec}
% \setotherlanguage{russian}
% \newfontfamily\cyrillicfont{Times New Roman}
% \newfontfamily\cyrillicfontsf{Times New Roman}[Script=Cyrillic]
% \newfontfamily\cyrillicfonttt{Times New Roman}[Script=Cyrillic]
%
% in the text use \begin{russian} ... \end{russian}
% LANGUAGE - END

% EMOJIS - BEGIN
% to use emoticons in your manuscript
% https://stackoverflow.com/questions/190145/how-to-insert-emoticons-in-latex/57076064
% using font Symbola, which has full support
% the font may be downloaded at:
% https://dn-works.com/ufas/
% add to preamble:
% \newfontfamily\Symbola{Symbola}
% in the text use:
% {\Symbola }
% EMOJIS - END

% LABEL REFERENCE TO DESCRIPTIVE LIST - BEGIN
% reference itens in a descriptive list using their labels instead of numbers
% insert the code below in the preambule:
%\makeatletter
%\let\orgdescriptionlabel\descriptionlabel
%\renewcommand*{\descriptionlabel}[1]{%
%  \let\orglabel\label
%  \let\label\@gobble
%  \phantomsection
%  \edef\@currentlabel{#1\unskip}%
%  \let\label\orglabel
%  \orgdescriptionlabel{#1}%
%}
%\makeatother
%
% in your document, use as illustraded here:
%\begin{description}
%  \item[first\label{itm1}] this is only an example;
%  % ...  add more items
%\end{description}
% LABEL REFERENCE TO DESCRIPTIVE LIST - END


% add line numbers for submission
%\usepackage{lineno}
%\linenumbers

\begin{document}
\maketitle

\begin{polyabstract}
\begin{abstract}
Os transtornos alimentares afetam até 86\% das crianças com Transtorno do Espectro Autista (TEA). Embora as intervenções baseadas na Análise do Comportamento Aplicada (ABA) sejam eficazes, o acesso a elas é limitado, especialmente em países em desenvolvimento. O treinamento parental por meio de aplicativos surgiu como uma alternativa viável, mas a usabilidade desse tipo de recurso tem sido pouco investigada. Este estudo teve como objetivo avaliar a usabilidade, por especialistas, de um aplicativo de treinamento parental para o manejo da seletividade alimentar em crianças com TEA. O aplicativo “Nosso Diário TEA” foi estruturado dentro do modelo de Treinamento de Habilidades Comportamentais (BST) em 24 etapas. Sete especialistas em ABA participaram e responderam ao Questionário de Usabilidade para Smartphones (SURE) e a duas perguntas abertas, sendo que as análises incluíram estatísticas descritivas e análise temática categorial com base nas heurísticas de Nielsen. Cinco dos sete especialistas avaliaram o aplicativo no nível 80, o que representa forte concordância com os parâmetros de usabilidade. A análise qualitativa revelou satisfação com a acessibilidade da ferramenta, embora contenha algumas falhas identificadas nas heurísticas de Nielsen. De modo geral, o aplicativo “Nosso Diário TEA” demonstra um alto nível de conformidade com os critérios de usabilidade, apesar das inconsistências identificadas. Recomendações para melhorias foram feitas e serão utilizadas em pesquisas futuras com usuários finais.

\keywords{Usabilidade \sep Dispositivo móvel \sep Telemedicina \sep Treinamento parental \sep Autismo}
\end{abstract}

\begin{english}
\begin{abstract}
Eating disorders affect up to 86\% of children with Autism Spectrum Disorder (ASD). Although Applied Behavior Analysis (ABA) interventions are effective, access to them is limited, especially in developing countries. Parental training through apps has emerged as a viable alternative, but the usability of this type of resource has been little investigated. This study aimed to evaluate the usability by experts of a parental training app for managing food selectivity in children with ASD. The “Nosso Diário TEA” app was structured within the Behavior Skills Training (BST) model in 24 steps. Seven ABA experts participated and responded to the Smartphone Usability Questionnaire (SURE) and two open-ended questions, and the analyses included descriptive statistics and categorical thematic analysis based on Nielsen’s heuristics. Five of the seven experts rated the app at level 80, which represents strong agreement with usability parameters. Qualitative analysis revealed satisfaction with the tool's accessibility, although it contains some flaws identified in Nielsen's heuristics. Overall, the “Nosso Diário TEA” app demonstrates a high level of compliance with usability criteria, despite the inconsistencies. Recommendations for improving the app were made and will be used in future research with end users.

\keywords{Usability \sep Mobile device \sep Telemedicine \sep Parental training \sep Autism}
\end{abstract}
\end{english}
% if there is another abstract, insert it here using the same scheme
\end{polyabstract}

\section{Introdução}\label{sec-intro}
O Transtorno do Espectro Autista (TEA) é uma condição do neurodesenvolvimento reconhecida por prejuízos em duas áreas centrais: comunicação, interação social e comportamentos repetitivos e/ou estereotipados \cite{apa2023dsm5tr}. Além disso, a literatura indica uma correlação entre os sintomas do autismo e problemas de comportamento alimentar \cite{paula2020impacto}. Esses comportamentos abrangem recusa alimentar, um repertório alimentar restrito e/ou uma ingestão alimentar unidirecional, que são caracterizados na literatura como seletividade alimentar \cite{leader2020feeding}. Há uma estimativa de que 46\% a 86\% das crianças com TEA são acometidas por distúrbios alimentares, variando desde perturbações sensoriais (cor, textura ou odor) até comportamentos desadaptativos (agressão ou autolesão, birra, gritos, ânsia ou indução de vômito) \cite{bandini2017changes,suarez2015food}.

Procedimentos de base analítico-comportamental têm sido amplamente adotados nas intervenções para TEA devido à robustez e sustentação empírica da Análise do Comportamento Aplicada (ABA) nesse campo \cite{steinbrenner2020evidencebased,yu2020aba}). Apesar dos resultados promissores, a eficácia da intervenção pode ser impactada pela elevada carga horária semanal exigida, a depender do déficit de repertório do aprendiz \cite{casp2020practice}.

Uma alternativa acessível para intervenções baseadas em ABA é o treinamento parental \cite{dawsonsquibb2020parent,ijaz2021parent}. A literatura evidencia resultados favoráveis provenientes do treinamento parental direcionado ao TEA \cite{althoff2019parent,lee2021parentmediated,nevill2018metaanalysis,steinbrenner2020evidencebased}. Um método baseado em ABA, frequentemente utilizado para o treinamento parental, é o Behavior Skill Training (BST), que consiste no ensino de novas habilidades através de uma sequência composta por: instrução, modelo, ensaio e feedback \cite{gianoumis2012effects,sarokoff2004behavioralskills}. Apesar de sua efetividade, a implementação tradicional do BST pode limitar sua escalabilidade por depender fortemente da presença contínua de profissionais especializados, favorecendo a investigação de formatos remotos e digitais \cite{erath2021training}. Nesse contexto, embora formatos remotos de BST tenham demonstrado efetividade semelhante ao presencial e ampliado o potencial de escalabilidade, também apresentam desafios técnicos e operacionais que impactam sua implementação em contextos aplicados \cite{lerman2020remote}. Além disso, existe uma escassez de estudos no que se refere ao treinamento parental baseado em BST voltado para a seletividade alimentar em crianças com TEA \cite{correaferreira2025treinamento}.

Embora escassos e com limitações amostrais, ensaios clínicos com estratégias de telessaúde e videomodelação direcionadas a pais e cuidadores têm demonstrado resultados positivos. \textcite{bloomfield2019food} utilizaram treinamento parental por teleconsulta baseado em BST para o manejo da seletividade alimentar de uma criança de 8 anos, observando aumento na aceitação de novos alimentos e alta aceitabilidade da intervenção. De forma semelhante, \textcite{gerow2021telehealth} treinaram pais via coaching telefônico como intervencionistas de crianças de 5 a 9 anos, verificando implementação com fidelidade e aumento das habilidades de vida diária em todos os participantes.

Todavia é necessária avaliação da usabilidade dos dispositivos tecnológicos antes da sua implementação prática para garantir que sejam seguros, eficazes e acessíveis aos usuários finais \cite{mclaughlin2020evaluating}. A mensuração da usabilidade é responsável por investigar a qualidade intrínseca de uma interface, além de avaliar a eficiência, eficácia e satisfação do utilizador com o uso de uma determinada tecnologia, inserida em um determinado ambiente para a realização de objetivos específicos \cite{cybis2015ergonomia,nielsen2014usabilidade}.

A avaliação da usabilidade está em consonância com uma das normas técnicas da International Organization for Standardization (ISO), que enfatiza a importância de desenvolver um sistema com eficácia, eficiência e satisfação para usuários específicos utilizarem a tecnologia em contextos específicos (ISO 9241-11, 2011). Nesse sentido, foram conduzidos ensaios com essa finalidade. \textcite{bonnot2021mobile} avaliaram a usabilidade e utilidade de um aplicativo de empoderamento parental, enquanto \textcite{parsons2019toby} exploraram a adequação de uma ferramenta pela perspectiva dos pais, incluindo métricas quantitativas de usabilidade. \textcite{mruzek2019toilettraining}, por sua vez, focaram na usabilidade e eficácia de um sistema para enurese. Entretanto, a revisão de escopo de \textcite{bharat2023mhealth} indica que a maioria dos aplicativos para TEA carece de avaliações de usabilidade antes de serem disponibilizados. Além disso, nenhum desses ensaios teve como objetivo principal o manejo da seletividade alimentar, embora alguns coletem dados sobre comportamento alimentar ou relatem desfechos secundários na rotina alimentar. Desse modo, existe uma escassez de evidências para a avaliação da usabilidade de aplicativos voltados ao treinamento parental para a seletividade alimentar no autismo, apesar do reconhecimento de que a melhoria da usabilidade pode resultar em maiores benefícios em termos de qualidade de vida para essa população \cite{mator2021usability}.

É importante destacar que a avaliação da usabilidade de uma nova tecnologia não deve ser realizada apenas pelo usuário final \cite{zulzalil2023checklistframework}. Especificamente com tecnologias voltadas ao público com TEA, é recomendada a avaliação inicial por especialistas, assegurando que o produto ou serviço ofereça uma experiência mais positiva e alinhada às suas particularidades \cite{valencia2022userexperience}. Desse modo, o objetivo principal desta pesquisa é avaliar a usabilidade por experts sobre um aplicativo de treinamento parental do manejo de seletividade alimentar de crianças com TEA.

\section{Método}\label{sec-normas}
Essa pesquisa é um recorte de um projeto maior, intitulado “Capacitação Familiar Digital: Manejo de Seletividade Alimentar em Crianças com Transtorno do Espectro Autista”, e visa promover o aprimoramento e continuidade de uma linha de pesquisa que vem sendo desenvolvida desde 2013, envolvendo diferentes universidades do Brasil \cite{bagaiolo2018capacitacao,bagaiolo2019community}. O projeto foi aprovado pelo Comitê de Ética em Pesquisa com Seres Humanos da Universidade do Vale do Itajaí (UNIVALI) sob o parecer nº 7.726.954 em 25 de julho de 2025.

\subsection{Aplicativo "Nosso Diário TEA"}\label{sec-conduta}
O aplicativo “Nosso Diário TEA”, disponível na plataforma do Google Play©, foi desenvolvido pelos pesquisadores desse estudo que contam com uma equipe composta por doutores em psicologia, estudantes de psicologia em nível de graduação e pós-graduação com experiência em ABA, bem como com a colaboração de uma nutricionista. A tecnologia já havia sido utilizada anteriormente para o treinamento parental da comunicação funcional e manejo de comportamentos disruptivos. Nesta fase do projeto, o aplicativo foi reestruturado para fornecer o treinamento parental em relação à seletividade alimentar.

Foram utilizados recursos de Inteligência Artificial para revisar a linguagem dos textos e permitir maior compreensão dos conteúdos do aplicativo. Seguiu-se as recomendações do Publication Manual of the American Psychological Association. A estrutura do treinamento foi realizada com base nas melhores evidências disponíveis na literatura sobre treinamento parental por meio de aplicativos \cite{correaferreira2025treinamento,kral2023mobilehealth}. Após o desenvolvimento do conteúdo, foi contratado um cartunista para a elaboração dos desenhos que compõem uma narrativa gráfica curta. Em seguida, uma equipe de atores e de produção audiovisual foi responsável pela gravação e edição dos vídeos. Posteriormente, uma equipe de Tecnologia da Informação realizou a programação do aplicativo e o disponibilizou para download na plataforma digital.

A tecnologia possui dois módulos. O primeiro é dividido em 10 partes, destinadas à obtenção da linha de base do repertório dos pais, através do ensino de conceitos básicos da ABA. O segundo módulo é composto por 14 partes voltadas à aproximação dos alimentos, totalizando 24 etapas. Os componentes de treinamento para os pais estão descritos na \Cref{tab1}.


\begin{table}[h!]
\begin{threeparttable}
\caption{Etapas do treinamento.}
\label{tab1}
\begin{tabular}{p{2cm}p{11cm}}
\toprule
Etapas & Descrição \\
\midrule

\multicolumn{2}{l}{Módulo 1 -- Linha de Base} \\
\hline
1 & O que é seletividade alimentar \\
2 & Escolha do alimento \\
3 & O que é comportamento respondente \\
4 & O que é comportamento operante \\
5 & Tríplice contingência \\
6 & Reforço positivo \\
7 & Reforço negativo \\
8 & O que é extinção \\
9 & Punição \\
10 & Reforço diferencial \\
\midrule

\multicolumn{2}{l}{Módulo 2} \\
\midrule
1 & Escolha de apenas um alimento \\
2 & Passo 1: tolerar o alimento \\
3 & Passo 2: visualizar o alimento \\
4 & Passo 3.1: pré-cheirar \\
5 & Passo 3.2: cheirar o alimento \\
6 & Passo 4: tocar o alimento \\
7 & Passo 5: interagir com o alimento \\
8 & Passo 6: tocar o alimento na boca/lábios \\
9 & Passo 7.1: videomodeling para mostrar para a criança \\
10 & Passo 7.2: tocar o alimento na língua/lamber \\
11 & Passo 8: morder o alimento \\
12 & Passo 9: mastigar o alimento \\
13 & Passo 10: comer o alimento \\
14 & Material extra: como manter o que foi aprendido \\
\bottomrule

\end{tabular}

\vspace{0.3cm}
 \source{Elaboração dos autores.}
\end{threeparttable}
\end{table}

Em consonância com os princípios do Behavioral Skills Training (BST), cada etapa do aplicativo foi estruturada com quatro componentes principais: (1) instrução didática, apresentada em formato de texto explicativo; (2) modelação, por meio de narrativa gráfica e vídeo demonstrativo; (3) ensaio, operacionalizado por meio de uma questão de múltipla escolha relacionada ao conteúdo apresentado; e (4) feedback de desempenho, fornecido automaticamente pelo aplicativo após cada resposta. O feedback consistia na indicação de acerto ou erro. Cada etapa permitia até três tentativas de resposta. A cada nova tentativa o grau de complexidade da questão era reduzido, com o objetivo de favorecer a aprendizagem do conteúdo. Para avançar para a etapa seguinte, o participante precisava responder corretamente à questão proposta. Exemplos da interface do aplicativo podem ser encontrados na \Cref{fig-1}.

\begin{figure}[h!]
\centering
\begin{minipage}{.95\textwidth}
\includegraphics[width=\textwidth]{Fig1.png}
\caption{Capturas de tela do aplicativo “Nosso Diário TEA”..}
\label{fig-1}
\source{Elaboração própria.}
\notes{A primeira parte da imagem refere-se à etapa de instrução, a segunda à representação gráfica, a terceira refere-se ao vídeo e a quarta refere-se à última etapa, questionário/prova.}
\end{minipage}
\end{figure}

\subsection{Procedimentos}\label{sec-fmt-manuscrito}
Os convites aos participantes foram compartilhados junto ao link de inscrição através das redes sociais (e.g. Instagram e Whatsapp) em março de 2025. No período de 50 dias, realizaram-se três divulgações da pesquisa e as postagens foram colocadas como públicas, permitindo a qualquer pessoa compartilhar o convite. Ao clicar no link, os participantes foram direcionados ao Termo de Consentimento Livre e Esclarecido (TCLE). Logo, os que aceitaram o TCLE informaram seus dados para contato. Posteriormente, os pesquisadores contataram os participantes via Whatsapp, e enviaram um link para preenchimento do questionário sociodemográfico. Para ser incluído, o participante precisava ser um profissional da saúde, possuir pós-graduação, especialização lato sensu em ABA, ter no mínimo um ano de experiência e ter acesso a um aparelho Android. Foram excluídos profissionais que atuaram diretamente no projeto, que estavam vinculados às instituições participantes da pesquisa.

Os participantes que atenderam aos critérios de inclusão receberam um link para baixar o aplicativo no Google Play©, junto a um login e senha para acessá-lo. Inicialmente, foi estabelecido um prazo de 10 dias para acessar o aplicativo e navegar pela intervenção. Durante esse período, os pesquisadores ficaram disponíveis para sanar dúvidas, além de enviaram mensagens com lembretes e incentivos verbais para o acesso.

Após finalizar o prazo, foi enviado um terceiro link com o questionário Smartphone Usability Questionaire (SURE) e duas perguntas abertas. O instrumento SURE é composto por 31 itens em uma Escala Likert de 5 pontos. O instrumento avalia diversos aspectos da usabilidade e da satisfação dos usuários em relação aos aplicativos de smartphones através de uma escala construída a partir da Teoria de Resposta ao Item \cite{wangenheim2014sure}. À medida que o escore total é obtido, classificam-se os resultados nas faixas de distribuição 30, 40, 50, 70 e 80 \cite{wangenheim2014sure}. Ao final, os participantes tiveram acesso a duas perguntas abertas elaboradas pelos pesquisadores: “comente sobre sua experiência utilizando o aplicativo”; e “descreva quais são suas sugestões para a melhora do aplicativo”.


\subsection{Participantes}\label{sec-formato}
Um total de 36 participantes aceitaram o TCLE. Destes, 10 completaram os questionários da pesquisa, porém, 3 foram excluídos por não possuírem pós-graduação em ABA. Os sete participantes incluídos relataram ter idade média de 28.57 anos (DP= 6.55), dos quais seis reportaram ser de gênero feminino. Todos os participantes eram residentes no Brasil, principalmente da região Sul (N=5). A maioria tinha graduação em psicologia (N=5). Além disso, 4 eram pós-graduados, 2 possuíam mestrado e 1 era especialista. Todos estavam formados há mais de dois anos, atuavam com ABA há mais de um ano e possuíam experiência com treinamento parental.


\section{Análise de dados}\label{sec-modelo}
Para a análise dos dados provenientes do SURE foi utilizado a análise estatística descritiva \cite{dancey2019estatistica}. Com relação às duas perguntas abertas, foi empregada a análise temática categorial de \textcite{bardin2011analise}, em categorias a priori a partir das dimensões das heurísticas de Nielsen \cite{nielsen1995heuristics}. As 10 heurísticas e suas definições são apresentadas na \Cref{tab2}.

\begin{table}[h!]
\begin{threeparttable}
\caption{Definições das heurísticas de Nielsen.}
\label{tab2}
\begin{tabular}{p{3.5cm}p{10cm}}
\toprule
Heurística & Descrição \\
\midrule

1. Visibilidade do status do sistema; \textit{feedback} &
Fornece aos usuários \textit{feedback} rápido e relevante para que sempre entendam seu status atual, ações disponíveis e como proceder. \\

2. Correspondência entre o sistema e o mundo real &
Utilize termos, frases e conceitos familiares alinhados à experiência dos usuários, para apresentar informações de forma natural e lógica. \\

3. Controle e liberdade do usuário &
Permite que os usuários controlem a navegação e o ritmo. Fornece maneiras óbvias de sair de caminhos indesejados e suporte ações de desfazer/refazer. \\

4. Consistência e padrões &
Garante consistência na linguagem, significado e comportamento em toda a interface, de modo a alinhar-se com as expectativas dos usuários e as normas padrão da plataforma. \\

5. Prevenção de erros &
Previne erros por meio de um design cuidadoso. Mesmo quando ocorrerem erros, o design deve permitir recuperação eficaz ou minimizar o impacto. \\

6. Reconhecimento em vez de recordação &
Exibe claramente objetos, opções e ações. Os usuários não devem precisar lembrar informações de diferentes partes do sistema; instruções essenciais devem estar facilmente acessíveis. \\

7. Flexibilidade e eficiência de uso &
Permite que os usuários adaptem o sistema às suas preferências e habilidades, incluindo recursos como atalhos para eficiência. \\

8. Design estético e minimalista &
Evite informações desnecessárias ou distrativas que concorram com o conteúdo relevante. Mantenha elementos visuais como animações mínimos e com propósito. \\

9. Reconhecer, diagnosticar e recuperar-se de erros &
Expresse mensagens de erro de forma clara, usando linguagem simples. Indique o que deu errado e ofereça orientação para resolver o problema. \\

10. Ajuda e documentação &
Priorize sistemas que sejam intuitivos e exijam documentação mínima. Se a documentação for necessária, assegure-se de que seja breve, orientada a tarefas e fácil de navegar. \\

\bottomrule
\end{tabular}

\vspace{0.3cm}
\notes{As dez heurísticas mencionadas acima podem ser encontradas em \textcite{nielsen1995heuristics}}.
\source{Elaboração própria.}
\end{threeparttable}
\end{table}



\section{Resultados}\label{sec-organizacao}
Nesta seção, serão apresentadas as análises estatísticas do SURE e análise de conteúdo das questões abertas respondidas pelos experts sobre a usabilidade. Além disso, com base nas heurísticas identificadas, serão levantadas algumas recomendações para o aprimoramento da tecnologia.

\subsection{Sure}\label{sec-organizacao-latex}
Na \Cref{tab3} apresenta-se a somatória da pontuação de cada participante e a classificação do nível de usabilidade obtido através do SURE. A média da somatória de pontuações de todos os participantes resultou em 88.57 (DP=15.65). Os dados obtidos demonstram que 5 participantes classificaram a ferramenta com nível de usabilidade 80, representando que concordam totalmente sobre a ajuda e o apoio do aplicativo para o alcance do seu objetivo. No entanto, um participante se enquadra no nível 70 de usabilidade, o que evidencia uma concordância forte sobre as afirmações e indicam que o aplicativo possui design e textos adequados, cumpre com as necessidades do consumo e recomendaria esse aplicativo para outros usuários, entre outras considerações. Por outro lado, o participante 6 apresentou classificação de usabilidade no nível 50, o que sugere uma percepção mais moderada em relação à experiência de uso, indicando a presença de aspectos que podem ser aprimorados na ferramenta.

\begin{table}[h!]
\begin{threeparttable}
\caption{Somatória de pontuações no SURE e classificação do nível de usabilidade.}
\label{tab3}
\begin{tabular}{ccc}
\toprule
Participante &
Somatória de pontuação no SURE &
Classificação do nível de usabilidade \\
\midrule

01 & 93  & Nível 80 \\
02 & 94  & Nível 80 \\
03 & 89  & Nível 80 \\
04 & 113 & Nível 80 \\
05 & 78  & Nível 70 \\
06 & 62  & Nível 50 \\
07 & 91  & Nível 80 \\

\bottomrule
\end{tabular}
\source{Elaboração dos autores.}
\end{threeparttable}
\end{table}

\subsection{Heurísticas de Nielsen}\label{sec-titulo}
As respostas dos especialistas às perguntas abertas permitiram identificar problemas e sugestões de usabilidade com base nas heurísticas de \textcite{nielsen1995heuristics}, resultando em recomendações para aprimorar o aplicativo (ver \Cref{tab4}).

\begin{longtable}{p{2.5cm}p{5cm}p{6cm}}
\caption{Categorias de análise dos problemas de usabilidade baseadas nas heurísticas de Nielsen.}
\label{tab4} \\
\toprule
Heurísticas de Nielsen &
Problemas de Usabilidade &
Soluções Recomendadas \\
\midrule
\endfirsthead

\midrule
Heurísticas de Nielsen&
Problemas de Usabilidade&
Soluções Recomendadas\\
\endhead

1. Visibilidade do status do sistema &
E2: “Não se sabe se uma atividade continua ou já terminou.”; E5: “Senti falta de uma notificação ou aviso claro ao final das atividades ou módulos — acabei de não perceber que havia conclusão do conteúdo proposto.” &
Incluir mensagens visuais e auditivas claras ao final de cada atividade/módulo (e.g., “Parabéns, você concluiu o primeiro módulo”); estabelecer sistema de recompensa gamificada ao concluir o módulo e também ao concluir a tarefa; substituir barra de progresso por uma trilha de visualização, na qual todos os módulos são visíveis e liberados à medida que o usuário avança; exibir um vídeo inicial apresentando toda a jornada do app, módulos e formato de tarefas logo na tela inicial. \\
\hline

2. Correspondência entre o sistema e o mundo real &
E6: “Pensando em aplicabilidade, sobretudo para pais, achei um pouco extenso, o que talvez aumenta o custo da resposta para essas pessoas [...]” &
Adequar linguagem acessível, com metáforas familiares e exemplos cotidianos para pais; adequar linguagem científica à comunicação direta para que o usuário reconheça termos ao invés de aprender novos. \\
\hline

5. Prevenção de erros &
E2: “O aplicativo tem problemas para abrir, trava e é preciso abrir e fechar novamente [...]. Há momentos em que o aplicativo retorna ao mesmo módulo [...]. Não é possível inserir dados no diário.”; E7: “Os vídeos demoram para carregar.” &
Otimizar o desempenho técnico do app, com testes de estabilidade e menor tempo de carregamento de mídia; implementar salvamento automático de progresso e correção de bugs na navegação entre módulos; corrigir o funcionamento do diário com testes de usabilidade em dispositivos diversos; reformular o diário, em formato gamificado, apenas com seleção do cumprimento da tarefa; estabelecer esquema de recompensa pelo uso do diário ou inserção de fotos. \\
\hline

7. Flexibilidade e eficiência de uso &
E1: “Colocar todos os módulos para que a pessoa saiba o que vai fazer…”; E2: “Talvez você possa esclarecer que eles podem assistir aos vídeos ou só às narrativas gráficas curtas…” &
Tornar a trilha de conteúdo visível desde o início, com opção de navegação livre por módulo (e.g., ver primeiro apenas vídeos ao invés de iniciar com instrução); introduzir interação gamificada para manter a simplicidade da interface, com botões intuitivos. \\
\hline

8. Estética e design minimalista &
E5: “Os vídeos, embora informativos, tornaram-se cansativos em certos momentos por serem longos e, em alguns casos, repetitivos [...]. As tirinhas, que poderiam ser um recurso mais leve e reflexivo, às vezes rasas, com pouco aprofundamento no conteúdo – apenas uma imagem, sem uma contextualização maior.”; E3: “Textos bem pequenos, algumas pessoas podem ter dificuldade para ler…” &
Reduzir o tempo dos vídeos, tornando-os mais dinâmicos; introduzir narração nos vídeos; ao final de cada vídeo, apresentar tópicos-chave apresentados, associando comunicação visual à comunicação escrita; reestruturar as narrativas gráficas curtas para histórias em quadrinhos com contexto textual e narração breve; reformular os textos para que sejam curtos e claros; aprimorar a personalização do app para aumentar tamanho da fonte e contraste claro e escuro; personalizar a coloração do app. \\
\hline

10. Ajuda e documentação &
E5: “Ausência de mensagens que ampararam o usuário em caso de travamentos.” &
Incluir mensagens de erro informativas e sugestões de solução (e.g., “Algo deu errado, tente reiniciar. Se persistir, clique aqui para ajuda”); criar uma seção de FAQ e tutoriais interativos no menu principal. \\
\hline
\source{Elaboração própria.}
\end{longtable}

A heurística 1 (visibilidade do status do sistema) destaca a importância de manter o usuário informado, por meio de feedbacks claros e em tempo adequado. No entanto, os especialistas identificaram falhas nesse aspecto no aplicativo avaliado. As narrativas dos usuários demonstram que eles enfrentaram dificuldades para compreender em que etapa estavam dentro do aplicativo. A ausência de marcadores de progresso, feedbacks de conclusão ou visualização da trilha completa do percurso comprometeu a percepção de avanço e continuidade. As sugestões apresentadas indicam uma demanda por maior previsibilidade e clareza no uso da plataforma, o que impacta diretamente a autonomia e o engajamento dos usuários durante a experiência de aprendizagem.

A heurística 2 (Correspondência entre o Sistema e o Mundo Real) orienta que o sistema utilize uma linguagem familiar ao usuário, com termos e elementos visuais que remetem ao seu cotidiano, facilitando a compreensão e reduzindo a carga cognitiva. No estudo, há divergências de opiniões com relação a este aspecto. Por um lado, os especialistas destacaram que o formato do aplicativo poderá provocar aumento da carga cognitiva dos usuários (\Cref{tab4}), enquanto outros experts ressaltaram aspectos positivos nesse quesito, principalmente acessibilidade, a utilização de recursos visuais com exemplos cotidianos e o respeito ao modelo teórico da ABA, conforme exemplifica o participante E4: “Gostei da experiência. Pela facilidade de acesso e base científica de condução, acredito que muitas famílias conseguirão utilizá-lo com facilidade”.

A heurística 5 (prevenção de erros) recomenda que o sistema seja projetado para prevenir erros, guiando o usuário com informações claras e seguras. Em aplicativos com sequência de atividades, como o avaliado, isso é especialmente importante. No entanto, os especialistas identificaram falhas que prejudicam a experiência do usuário. Os experts apontaram erros operacionais recorrentes, como travamentos e falhas no carregamento, que afetam diretamente a confiança e a fluidez do uso. Tais falhas comprometem não apenas a usabilidade, mas poderá interferir no engajamento e na continuidade do processo de aprendizagem dos pais. A prevenção de erros demonstra ser uma heurística crítica e ainda fragilizada no aplicativo analisado.

A heurística 7 (flexibilidade e eficiência de uso) defende que o sistema seja flexível, atendendo tanto iniciantes quanto usuários experientes, por meio de instruções claras, atalhos e personalização. No aplicativo avaliado, os especialistas consideraram a importância de oferecer uma visão geral do percurso de aprendizagem, favorecendo o planejamento e a autonomia do usuário e a necessidade de maior clareza e flexibilidade na escolha dos conteúdos, permitindo que cada usuário personalize sua experiência conforme sua preferência ou necessidade. Cabe ressaltar que em outras heurísticas os experts apontaram dificuldades técnicas e de navegação – como travamentos, retornos inesperados e lentidão – que comprometem a presente heurística de eficiência do uso. Portanto, os aspectos de flexibilidade e eficiência ainda exigem aprimoramentos que atendam tanto usuários iniciantes quanto mais experientes. Apesar disso, vários avaliadores reconheceram a sua facilidade de acesso, como exemplificado nas seguintes narrativas: E1: “Adorei utilizá-lo, é fácil e funcional.”; E3: “Foi bem fácil e tranquilo para mim de usar, me auxiliou muito no meu dia a dia!”; E4: “Facilidade de acesso [...], muitas famílias conseguirão utilizá-lo com facilidade.”; E7: “O aplicativo é de fácil acesso e entendimento”.

A heurística 8 (Design Estético e Minimalista) reforça que o design do sistema seja simples, com elementos visuais que não distraiam nem sobrecarreguem o usuário com informações irrelevantes. A estética deve ser funcional, facilitando a navegação, a leitura e o foco nos conteúdos principais. A avaliação dos experts indicou que o uso excessivo de textos, vídeos longos e imagens pode demandar mais tempo e esforço cognitivo, dificultando o engajamento e o aprendizado dos usuários.

A heurística 10 (Ajuda ao Usuário e Feedback) destaca a importância de oferecer suporte acessível e alinhado às possíveis dificuldades do usuário. A ajuda precisa ser fácil de encontrar, integrada à experiência de uso e capaz de responder claramente às dúvidas ou problemas que o usuário enfrentará no momento. A avaliação dos experts evidencia uma fragilidade no sistema de suporte diante de situações inesperadas, como falhas técnicas ou travamentos, o que pode comprometer a fluidez da navegação, gerar insegurança no usuário e afetar negativamente sua autonomia e continuidade no uso da plataforma.

Dentre as heurísticas de usabilidade propostas por Nielsen, quatro não apresentaram correspondência direta com as falas dos especialistas dessa pesquisa: heurística 3 (Controle e Liberdade do usuário), heurística 4 (Consistência e padrões), heurística 6 (Reconhecimento em vez de recordação) e heurística 9 (Reconhecer, diagnosticar e recuperar-se de erros).

A avaliação heurística de Nielsen permitiu o levantamento de sugestões recomendadas em resposta aos problemas de usabilidade identificados pelos especialistas para o aprimoramento do aplicativo. Dentre as principais recomendações destacam-se a necessidade de melhoria da visibilidade do status do sistema por meio de mensagens claras de conclusão, recompensas gamificadas e uma trilha visual de progresso acessível. Também é recomendado a adequação da linguagem ao público-alvo, prevenção de erros com melhorias técnicas e estabilização do aplicativo, promover o uso do aplicativo de maneira flexível através da navegação livre, bem como priorizar um design mais minimalista com conteúdos mais sucintos e recursos multimodais. Essas sugestões visam a correção das falhas e o aumento da satisfação da experiência do usuário para a utilização do aplicativo de forma mais intuitiva e engajadora.

\section{Discussão}\label{sec-autores}
Este é o primeiro estudo conhecido a avaliar a usabilidade sobre um aplicativo de treinamento parental do manejo de seletividade alimentar de crianças com TEA. Apenas dois aplicativos, TOBY e MITA, apresentam indicadores de usabilidade em aplicativos voltados aos pais de crianças com TEA, mas nenhum deles é específico sobre seletividade alimentar \cite{bharat2023mhealth}.

Os resultados obtidos com base nas respostas dos especialistas indicaram que a usabilidade do aplicativo Nosso Diário TEA foi, em geral, bem avaliada, principalmente considerando os dados do SURE, com cinco dos sete experts classificando a ferramenta no nível 80 de usabilidade, o que indica forte concordância com os critérios de utilidade, clareza e apoio ao alcance dos objetivos propostos pelo aplicativo. Os estudos que utilizaram o SURE para avaliar a usabilidade de um aplicativo e que obtiveram um nível maior ou igual a 70 indicam boa usabilidade para os objetivos do protótipo \cite{wangenheim2014sure,marques2020usability,silva2024qualiplus}. Embora não haja consenso sobre um número ideal de avaliadores para todos os testes de usabilidade, recomenda-se que estudos exploratórios de detecção de problemas contem com 5 a 10 participantes \cite{macefield2009participantsize}.

Todavia, um dos participantes dessa pesquisa classificou a usabilidade do aplicativo com nível inferior a 70, o que indica variações individuais na experiência de uso e que requerem maior aprofundamento com a contemplação de dados qualitativos sobre a percepção dos experts. A análise qualitativa em estudos sobre usabilidade é fundamental para compreender os aspectos subjetivos e contextuais que não são facilmente mensurados \cite{mator2021usability}. A análise qualitativa com base nas heurísticas de Nielsen permitiu aprofundar a compreensão sobre os aspectos específicos que impactam a experiência de uso. Os resultados revelaram que seis das dez heurísticas foram identificadas nesse estudo. Ressalta-se que não há necessidade de contemplar todas as heurísticas para avaliação da usabilidade de aplicativos na área da saúde, mas o seu conjunto deve ser suficientemente abrangente e relevante para o contexto específico de uso do aplicativo \cite{galavi2024heuristics}.

Dentre as heurísticas identificadas, destaca-se a heurística 1 (Visibilidade do status do sistema), os especialistas apontaram a ausência de indicadores de progresso e conclusões de módulos. Essa limitação impacta negativamente a autonomia e o engajamento dos usuários. De acordo com \textcite{bouraghi2022evaluating}, quando o status do sistema não é visível, os usuários tendem a cometer mais erros, apresentar menor produtividade e não manterem o engajamento ao longo do processo. Os autores sugerem o uso de elementos como barras de progresso e mensagens de conclusão para reforçar a clareza do percurso e motivar a continuidade, especialmente em sistemas com múltiplas etapas e objetivos educacionais. Neste sentido, percebe-se que por se tratar de um aplicativo formatado em etapas que exigem o cumprimento de um percurso de uso, este aspecto é fundamental para sua usabilidade. No estudo de \textcite{bonnot2021mobile}, com pais de crianças com TEA, é evidenciado que os usuários do aplicativo com maior taxa de resposta demonstraram maior interesse e satisfação com a tela de feedback, essa funcionalidade pode auxiliar no engajamento dos usuários sobre o seu desempenho no aplicativo.

A heurística 2 (Correspondência entre o sistema e o mundo real) uma das principais falhas nesse quesito refere-se ao aumento do custo de resposta dos usuários porque adota uma linguagem técnica. \textcite{stonbraker2018usability} destacam que a compreensão do conteúdo por públicos não especializados no tema exige uma adaptação do vocabulário, tornando o conteúdo familiar. A presença de termos técnicos da Análise do Comportamento Aplicada reforça a necessidade de adaptação da linguagem aos familiares, uma vez que o uso excessivo de jargões pode dificultar o engajamento dos cuidadores em consonância com a oitava dimensão da ABA, chamada de “compaixão”, que enfatiza a comunicação acessível e empática e centrada nas necessidades das famílias \cite{penney2023compassion}.

Os especialistas identificaram problemas técnicos como travamentos, falhas no carregamento de módulos e vídeos que são atribuídos a heurística 5 (Prevenção de erros). Essas dificuldades interrompem o fluxo da atividade, exigem reinício de etapas e comprometem a experiência do usuário. \textcite{bouraghi2022evaluating} apontam que falhas técnicas e longos tempos de resposta são fatores que comprometem o fluxo de aprendizagem e engajamento no contexto digital, podendo afetar tanto a adesão quanto a eficácia da intervenção. Essas interrupções, além de quebrarem a continuidade do processo de aprendizagem, podem gerar frustração e abandono do uso.

Embora os participantes tenham apontado que o aplicativo é fácil e acessível, também sugeriram melhorias relacionadas à personalização da experiência. Estes conteúdos se relacionam com a heurística 7 (Flexibilidade e eficiência de uso), ao indicar que, embora o aplicativo atenda bem usuários iniciantes, pode limitar a experiência de usuários que buscam autonomia e liberdade de escolha de como navegar, o que pode ser útil quando se fala de pais de crianças com TEA, que apresentam limitações expressivas no cotidiano. \textcite{bouraghi2022evaluating} destacam que a falta de flexibilidade e personalização impacta diretamente na produtividade dos usuários. Em aplicativos de saúde, em que o tempo e o engajamento são fatores críticos, oferecer opções de caminhos alternativos e permitir diferentes formas de interação com o conteúdo torna-se fundamental para ampliar a adesão e a eficácia da proposta. Além disso, no contexto geral da usabilidade para aplicativos na área da saúde, essa heurística é relevante para garantir que pais possam utilizar o aplicativo de forma eficaz e segura \cite{galavi2024heuristics}.

Com relação ao aspecto visual e funcional do aplicativo, os especialistas apontaram que vídeos longos, narrativas gráficas curtas superficiais e excesso de texto comprometem a leveza da interface e aumentam o esforço cognitivo, o que está relacionado a violações da heurística 8 (Design estético e minimalista). \textcite{labrie2020adapting} mostram que ambientes digitais sobrecarregados com múltiplos elementos visuais geram desorganização e dificultam a compreensão, especialmente quando o conteúdo exige atenção prolongada. Nesse sentido, recomenda-se o uso de instruções mais objetivas, fragmentação de vídeos longos e que mantém o usuário engajado, além de adequação da densidade textual ao público-alvo. Essas estratégias visam otimizar a atenção, facilitar a compreensão e reduzir o tempo necessário para a execução das tarefas. Estes aspectos corroboram para que o uso do sistema seja articulado à rotina do usuário. No estudo sobre a avaliação da usabilidade do aplicativo Smartautism para pais de crianças com TEA, a limitação na qualidade estética da ferramenta pode ter sido uma das barreiras que dificultaram a adesão dos pais \cite{bonnot2021mobile}.

Os aspectos voltados às violações da heurística 10 (Ajuda ao usuário e feedback), são observados na fala dos especialistas quando há ausência de mensagens de erro ou orientações em momentos de falha que podem comprometer diretamente a experiência do usuário. \textcite{shen2024evaluating} argumentam que, para que os sistemas sejam eficientes, é necessário que ofereçam ajuda contextualizada e suporte imediato diante de dificuldades técnicas ou operacionais. A ausência de um canal de suporte ou de mensagens automáticas que orientem o usuário em caso de travamentos ou falhas pode gerar sensação de abandono e aumentar a frustração, especialmente em momentos em que o usuário está mais vulnerável ou sobrecarregado cognitivamente. \textcite{mruzek2019toilettraining} apontam que uma das melhorias necessárias no seu dispositivo, que contempla os pais para o treino à toalete de crianças com TEA, envolve o acesso eletrônico a manuais e suporte visual.



\section{Conclusão}\label{sec-resumo}
A avaliação da usabilidade de novas tecnologias digitais aplicadas aos pais de crianças com TEA é fundamental porque permite compreender com maior acurácia os parâmetros de eficiência, eficácia e satisfação. Nesse sentido, o presente estudo evidenciou que a usabilidade do aplicativo Nosso Diário TEA, avaliada por experts demonstra satisfação com acessibilidade e a maioria dos participantes atribuiu níveis elevados na sua classificação, o que evidencia uma forte concordância com os padrões de usabilidade. Embora falhas do aplicativo baseadas nas seguintes heurísticas de Nielsen foram identificadas na análise qualitativa: heurística 1 (Visibilidade do status do sistema), heurística 2 (Correspondência entre o sistema e o mundo real), heurística 5 (Prevenção de erros), heurística 7 (Flexibilidade e eficiência de uso), heurística 8 (Estética e design minimalista), heurística 10 (Ajuda e documentação). Foram apontadas recomendações para cada um desses problemas de usabilidade com o objetivo de aprimorar o aplicativo através de uma futura pesquisa de avaliação com usuários finais para garantir uma experiência mais satisfatória e intuitiva na utilização do aplicativo, bem como fortalecer seu processo de validação.

Dentre as principais limitações da pesquisa se refere a falta de benchmarks específicos para a avaliação da usabilidade de aplicativos sobre treinamento de pais de crianças com TEA e seletividade alimentar. Essa ausência de parâmetros de medida dificultaram a comparação dos resultados obtidos neste estudo com outras pesquisas, principalmente porque a literatura sobre o tema é escassa.

Outra limitação importante se refere ao fato de o aplicativo estar disponível apenas para dispositivos com sistema operacional Android, o que restringe o acesso de potenciais usuários que utilizam outras plataformas. Soma-se a isso o perfil dos participantes, composto exclusivamente por especialistas em Análise do Comportamento Aplicada, não incluindo profissionais da área de usabilidade ou design de interação, o que poderia ampliar a análise técnica da interface. Além disso, o número reduzido de participantes também pode ser considerado uma limitação, ainda que esteja em consonância com estudos de usabilidade conduzidos em fases iniciais de desenvolvimento de tecnologias em saúde. Observa-se ainda que os participantes estavam majoritariamente concentrados na região Sul do Brasil, não contemplando especialistas de outras regiões do país, onde a formação avançada em ABA também tem se expandido. Essa concentração regional pode limitar a diversidade de perspectivas culturais, formativas e contextuais sobre o uso da tecnologia. Como encaminhamento futuro, recomenda-se a replicação do estudo com um número maior de especialistas, incluindo profissionais de diferentes regiões do país e, preferencialmente, também especialistas em usabilidade. Além disso, uma etapa fundamental para o avanço da validação da tecnologia envolve a avaliação com os usuários finais, ou seja, pais e cuidadores de crianças com TEA, público-alvo direto da intervenção, permitindo analisar não apenas a usabilidade percebida, mas também a aplicabilidade prática e a compreensão do conteúdo no contexto cotidiano.

Em última análise, os resultados obtidos no estudo são promissores para a validação de uma tecnologia que poderá facilitar o acesso das famílias de crianças com TEA a uma intervenção parental sustentada por práticas baseadas em evidências da ciência ABA. Principalmente, ao se considerar que no contexto brasileiro os serviços especializados destinados a essa população são insuficientes diante da crescente demanda.

\begin{spacing}{0.9}
\printbibliography\label{sec-bib}
\end{spacing}
% if the text is not in Portuguese, it might be necessary to use the code below instead to print the correct ABNT abbreviations [s.n.], [s.l.]
%\begin{portuguese}
%\printbibliography[title={Bibliography}]
%\end{portuguese}


%full list: conceptualization,datacuration,formalanalysis,funding,investigation,methodology,projadm,resources,software,supervision,validation,visualization,writing,review
\begin{contributors}[sec-contributors]
\authorcontribution{João Rodrigo Maciel Portes}[conceptualization,datacuration,formalanalysis,projadm,investigation,methodology,supervision,visualization,writing,review]
\authorcontribution{Weliton da Silva}[conceptualization,datacuration,investigation,formalanalysis,visualization,writing]
\authorcontribution{João Vitor Girardi Turqueti da Rocha}[conceptualization,datacuration,investigation,formalanalysis,visualization,writing]
\authorcontribution{Letícia Corrêa Ferreira}[conceptualization,datacuration,investigation,formalanalysis,visualization,writing]
\authorcontribution{Leila Bagaiolo}[conceptualization,methodology,formalanalysis,projadm,visualization,review]
\authorcontribution{Mauro Luís Vieira}[conceptualization,methodology,formalanalysis,supervision,projadm,visualization,review]
\end{contributors}
 

\begin{dataavailability}
\txtdataavailability{dataonly} % options: dataavailable, dataonly, databody, datanotav, nodata
\end{dataavailability}

\begin{funding}
Essa pesquisa recebeu apoio financeiro da Fundação de Amparo à Pesquisa e Inovação do Estado de Santa Catarina (FAPESC) através dos editais de Chamada Pública FAPESC nº 54/2022 e nº 20/2024.
\end{funding}


\end{document}

